\documentclass[12pt]{article}
\usepackage{amsmath,amssymb,amsthm,fullpage,microtype,mathpazo,graphicx,wrapfig}
\usepackage[T1]{fontenc}

\include{hyper}

\renewcommand{\r}{{\bf r}}
\newcommand{\dery}{\frac{dx}{dt}}
\renewcommand{\k}{{\bf k}}
\newcommand{\be}{\begin{equation}}
\newcommand{\ee}{\end{equation}}
\newcommand{\ns}{\vspace{3mm}}
\newcommand{\Dth}{\Delta \theta}
\newcommand{\p}{\partial}

\setlength{\fboxsep}{8pt}

\newcommand{\exer}[2]{
\begin{center}
  \vspace{2mm}
\framebox{\parbox{15cm}{
{
\textbf{\large Exercise #1}\vspace{1mm} \\
#2
}}}
  \vspace{2mm}
\end{center}
}

\usepackage{hyperref,color}

\definecolor{webred}{rgb}{0.7,0,0.1}
\definecolor{webgreen}{rgb}{0,.35,0}
\definecolor{webbrown}{rgb}{.6,0,0}
\definecolor{RoyalBlue}{rgb}{0,0,0.9}
\definecolor{mywhite}{rgb}{1.0,1.0,1.0}
\newcommand{\white}[1]{\textcolor{mywhite}{#1}}

\hypersetup{%
   colorlinks=true, linktocpage=true, pdfstartpage=3, pdfstartview=FitV,
   breaklinks=true, pdfpagemode=UseNone, pageanchor=true, pdfpagemode=UseOutlines,
   plainpages=false, bookmarksnumbered, bookmarksopen=true, bookmarksopenlevel=1,
   hypertexnames=true, pdfhighlight=/O,
   urlcolor=webbrown, linkcolor=RoyalBlue, citecolor=webgreen,
   pdfauthor={Chris H. Rycroft},
   pdfsubject={Harvard AM50 (Spring 2014)},
   pdfkeywords={},
   pdfcreator={pdfLaTeX},
   pdfproducer={LaTeX with hyperref}
}

\renewcommand*{\thefootnote}{\fnsymbol{footnote}}

\begin{document}
\section*{AM50: Rosalind Franklin and the discovery of the double helix\footnote{Based on notes from Professor Rafael Jaramillo, MIT.}}
\subsection*{Biography}
Franklin was born in 1920 to a well-off London banking family.  After primary
boarding school she attended St.~Paul's Girls' School, the sister campus to the
famous St.~Paul's Boys' school that had educated the cream of English society
since 1506. Franklin was a privileged child, active in her upper class London
Jewish society and school sporting clubs and spending weekends at the family
country estate. Her relatively high position in the English class system saw
her well during her time at Cambridge and in Paris, but later became a sticking
point in her relationships in the decidedly more working class environment at
King's College. To paraphrase from {\it My Fair Lady}, the moment and
Englishman speaks he makes some other Englishman despise him. This dynamic was
certainly present with Franklin at King's.

She entered Cambridge in 1938 and stayed on for her doctorate through the war.
Franklin was conflicted over how she could stay in academia while most of her
family was directly helping the war effort. This conflict was partially
resolved when she found a highly utilitarian Ph.D. project working for the
British Coal Utilisation Research Association on porosity in coal.

Franklin had always been more comfortable in continental society than among her
fellow English, and after the war she landed a dream job at the State Central
Chemical Services Laboratory in Paris. In a disarming letter to her French
friend and mentor Adrienne Weill she encapsulated what many of us feel upon
emerging from our narrowly defined Ph.D. projects:
\begin{quotation}
  {\it If ever you hear of anybody anxious for the services of a physical
  chemist who knows very little about physical chemistry but a lot about the
  holes in coal, please let me know.}
\end{quotation}
In Paris she learned the techniques of X-ray diffraction, and built an
international reputation for her continued work on the structure of coals.
Those were probably the happiest years of her life. The image on the screen was
taken in the Paris lab and shows Franklin carrying the ritual after-lunch
coffee in labware.

In 1951 Franklin returned to London to a job in the biophysics group at King's
College. Her return was motivated by persistent suggestions from her large,
tight-knit family that she should be closer to home.  But in every other way
she dreaded returning to England and trading the heady atmosphere of the Paris
Left Bank for dowdy King's College. Her two years at King's were highly
productive, but also turbulent and unhappy, and in 1953 she gladly left for
Birkbeck College. At Birkbeck she turned from DNA to the structure of viruses
and spent the most productive years of her life. 

Franklin died in 1958 at the age of 37 from ovarian cancer, five years before
the Nobel Prize that went to Watson, Crick, and Wilkins. At the time of her
death she was a Research Associate in Biophysics at Birkbeck College and was
busily engaged with Aaron Klug on work on the structure of viruses that was
later recognized in Klug's 1982 Nobel Prize. 

I've marked a few select historical events along the timeline of her life. The
Battle of Britain in 1940 took place while Franklin was preparing for her
undergraduate degree exam. Erwin Schrodinger's influential book
{\it What is Life?}~was published in 1944. In 1951 Linus Pauling published the alpha-helix protein
structure; from that time on it was said that helices were in the air. Two
years later, in 1953, the three papers describing the structure of DNA appeared
in Nature.

\subsection*{The Race}
The story of the double helix traces the relationships between a small number
of scientists at two institutions, King's College London (shown on the right)
and Cambridge University (and specifically the Cavendish Laboratory, shown on
the left).

After World War II, molecular biology was the next big thing in science. Erwin
Schrodinger's book {\it What is Life?}, published in 1944, stimulated interest
in biological problems in the physical science community and is often credited
with giving birth to the field of biophysics. One of the largest groups in
England was assembled at King's College London by John Randall, who was
considered somewhat of a hero for his invention of the cavity magnetron during
the war. Randall's deputy director of biophysics was Maurice Wilkins, who had
spent the war working on the Manhattan Project in California but soon after
read Schrodinger's book and switched into biophysics. At King's Wilkins was
working on X-ray diffraction studies of nucleic acid with the help of a Ph.D.
student named Ray Gosling. At the same time a biophysics group was being
assembled at the Cavendish Laboratory at Cambridge under the direction of
William Lawrence Bragg (the younger of the two Braggs and also Randall's Ph.D.
advisor). In 1951 a biochemist named Francis Crick joined a team working on
protein structure at the Cavendish.

In 1951 Randall hired Rosalind Franklin, a physical chemist with an
international reputation for her X-ray diffraction work on the structure of
coals. Rosalind was hired to study the structure of proteins in solution, but
as her start date approached Randall became increasingly interested in DNA. His
interest was driven in part by an extraordinary batch of samples that had been
produced by Rudolf Singer in Berne and that had been distributed freely to
Wilkins (and others) at a meeting in 1950. Wilkins and Gosling had already
shown that these samples gave excellent X-ray diffraction data. In a fateful
letter to Franklin, Randall assigned her to the DNA problem and gave her
supervision over Gosling. Wilkins, who was cut out of the DNA work, had not
been consulted.

Franklin had joined a well funded but poorly managed lab. Wilkins was badly
stung by being ousted from the work on DNA, and after starting off on this
wrong foot the relationship between Wilkins and Franklin only got worse. Their
personality differences compounded the awkward situation at work and soon each
was complaining about the other to their respective friends. Franklin was an
enthusiastic hostess and quickly gained a reputation for throwing excellent
dinner parties. The guest lists often included colleagues from King's, but
never Wilkins. Within a short time the rift between Franklin and Wilkins was
widely known within the biophysics community at King's and Cambridge, as was
the fact that Franklin was obtaining remarkable data. Word even travelled as
far as Pasadena. Linus Pauling bluntly asked Randall to see the King's data,
under the assumption that King's was sufficiently dysfunctional that they would
not be able to analyze their own results. Pauling was rebuffed, but the English
community was put on alert that Pauling was on the DNA trail.

Wilkins, for his part, began spending more time in Cambridge. He and Crick had
met during the war and had remained friends. Crick and his wife presided over a
lively social scene in Cambridge where, unlike in London, Wilkins was always
welcome to dinner. That same year saw the arrival in Cambridge of another man
with keen interest in DNA, James Watson. Watson and Crick got along
marvelously, and together they peppered Wilkins with questions about Franklin's
work on DNA. Late in 1951 Watson and Crick built their first DNA model, and
invited the King's group up for a review. Franklin wasted no time in making
clear that their model was patently incorrect.  Bragg was embarrassed for his
institution and ordered all DNA work at Cambridge to a halt. Bragg and Randall
quietly agreed that all DNA work would be left to King's, and Watson and Crick
were made to send their bits of metal and wire down to London.

After the agreement between the lab directors, the DNA work at King's continued
throughout 1952 with an extreme sense of urgency. Watson remained highly
curious, and was often seen hanging around King's. And Franklin knew that
Wilkins was in near-constant communication with Watson and Crick about her
work. A preprint containing a proposed DNA model from Pauling added to the
pressure; Pauling's model was wrong, but that only bought the English groups a
small amount of time. Adding to the sense of urgency, in the middle of 1952
Franklin decided that she would leave King's for a position at Birkbeck
College, making this her final year working on DNA. 

The critical encounter took place on January 30, 1953. Watson was on one of his
frequent trips to London, and ran into Wilkins in the hallway. After some time
spent listening to Wilkins complain (again) about his treatment at King's,
Watson was shown photo 51. As he described in The Double Helix, ``The instant I
saw the picture my mouth fell open and my pulse began to race.'' Over dinner
Watson pressed Wilkins for numbers to go with the diagram, and began sketching
models on the train ride home. The very next day Bragg was convinced to break
the agreement with Randall, allowing Watson and Crick to resume model building.
In February, Watson and Crick obtained a copy of an up-to-date report on the
King's work presented to the Medical Research Council, the primary funding
agency for biophysics in England. The copy was given to them by Max Perutz, a
colleague at Cambridge and a member of the Research Council committee. This
detailed report, combined with the information Watson had from photo 51, was
nearly all that Watson and Crick needed to complete their model, which was
finished on March 7.

By the middle of March 1953 the word at King's was that the Cavendish had
cracked the structure of DNA. Randall was furious, but for Franklin the news
was irrelevant. She was looking forward to her new position at Birkbeck, and
she was confident that the Cambridge model was just that---a model made of
scrap metal and cardboard that would rise or fall on the experimental evidence.
She was unaware that her best data had been central to the Cambridge discovery,
and with several papers in preparation she was confident that her contribution
to the field rested on solid ground. The ensuing rush to publish was
accompanied by negotiations between King's, Cambridge, and the editorial
offices at Nature. The importance of the King's data to the Cambridge model was
not openly acknowledged, and it's unclear who at the time knew the full story.
Probably not the editors at Nature. Of the trio of papers on DNA that appeared
on April 25, the classic by Watson and Crick with its incomplete acknowledgment
came first, and the Franklin and Gosling paper containing photo 51 was placed
last.

One of the interesting ``What Ifs'' in this story is whether Franklin would have
been first to the DNA structure, had the Cambridge group not had such direct
access to her unpublished data. At the time of the discovery she was busily
engaged in Patterson analysis, a highly quantitative but time-consuming
technique that likely would have led to the double helix. Besides the
diffraction analysis, many of the stereochemical ideas that guided Watson and
Crick were also understood at King's. But as we all know, the seemingly short
distance between having 99\% of the results and telling a complete physical
story can in fact be quite long. Speaking after her death, and after carefully
reviewing her notebooks, her good friend, collaborator, and beneficiary Aaron
Klug concluded that:

\begin{quotation}
  {\it She needed a collaborator, and she didn't have one. Somebody to break
  the pattern of her thinking, to show her what was right in front of her, to
  push her up and over.}
\end{quotation}


\end{document}
